% Created 2026-03-26 qui 18:05
% Intended LaTeX compiler: lualatex
\documentclass[12pt,a4paper]{article}
\usepackage{fgv-paper}
\usepapercover
\institution{Faculdade Getúlio Vargas}
\programname{Mestrado Profissional em Políticas Públicas e Governo}
\coursename{}
\disciplinename{Política Brasileira Contemporânea}
\professorname{Raphael Amorim Machado}
\cityname{Brasília}
\papertype{Artigo acadêmico de análise bibliográfica e argumentativa}
\covernote{A estabilidade democrática não elimina o conflito; apenas lhe confere forma institucional. Compreender a ruptura exige examinar as mediações políticas que sustentam, e por vezes deixam de sustentar, a governabilidade.}
\author{Gustavo M. Mendes de Tarso}
\date{$\backslash$\today}
\title{Governabilidade, representação e ruptura presidencial na democracia brasileira pós-1988}
\hypersetup{
 pdfauthor={Gustavo M. Mendes de Tarso},
 pdftitle={Governabilidade, representação e ruptura presidencial na democracia brasileira pós-1988},
 pdfkeywords={},
 pdfsubject={},
 pdfcreator={Emacs 28.2 (Org mode 9.5.5)}, 
 pdflang={Pt_Br}}
\usepackage[style=backend=biber,style=apa,sortcites=true,sorting=nyt,giveninits=true,maxcitenames=2,maxbibnames=20,uniquelist=minyear]{biblatex}
\addbibresource{/home/gustavodetarso/Documentos/mppg/disciplinas/politica_brasileira_contemporanea/paper/paper_politica_brasileira_contemporanea/paper_politica_brasileira_contemporanea.bib}
\begin{document}

\makemytitle

\begin{abstract}
Este paper analisa a democracia brasileira pós-1988 a partir da relação entre governabilidade, presidencialismo de coalizão, sistema partidário e tensões de representação, tendo o impeachment de 2016 como aplicação analítica central. O problema de pesquisa consiste em explicar por que um arranjo institucional frequentemente associado à produção de estabilidade governamental também pode desembocar em ruptura presidencial sem que isso implique, necessariamente, sua inviabilidade estrutural. A pergunta que orienta o trabalho é a seguinte: em que medida o impeachment de 2016 pode ser explicado menos como evidência do fracasso intrínseco do presidencialismo de coalizão e mais como resultado da interação entre erosão da base partidária de governo, dinâmica competitiva do sistema partidário e tensões de representação? O objetivo é analisar em que medida estabilidade e crise presidencial podem ser mais bem compreendidas por essa interação, em vez de por explicações monocausais.

A hipótese sustenta que o presidencialismo de coalizão, tal como operou na democracia brasileira, demonstrou capacidade de coordenar maiorias e produzir governabilidade, como argumentam \textcite{palermo2000_como_se} e \textcite{limongi2017_a_crise}. A ruptura ocorre quando essa capacidade institucional deixa de estar lastreada por coalizões partidárias operantes e por mediações representativas minimamente estáveis. Nesses contextos, a recomposição adversa das maiorias legislativas, combinada à intensificação da polarização social e à deslegitimação dos partidos, torna o impeachment uma via institucionalmente disponível para converter impasses políticos em afastamento presidencial, em linha com a abordagem comparativa de \textcite{perez_lian2016_explaining_military}.

Metodologicamente, o paper adota abordagem bibliográfica e argumentativa, articulando textos-base da disciplina sobre instituições, partidos, representação, protesto e crise política. O argumento central é que o caso de 2016 não invalida, por si só, o diagnóstico de que o sistema político brasileiro foi capaz de produzir coordenação decisória; ao contrário, ele revela os limites condicionais desse padrão de governabilidade quando se alteram simultaneamente os incentivos da competição partidária e os vínculos de representação. A relevância do trabalho reside em oferecer uma interpretação que conecta estrutura institucional e dinâmica política, evitando tanto a leitura institucionalista simplificadora quanto a narrativa episódica do impeachment.
\end{abstract}

\section{Introdução}
\label{sec:orga5ead32}
A democracia brasileira pós-1988 foi frequentemente interpretada a partir de uma tensão entre desenho institucional complexo e capacidade efetiva de governo. Presidencialismo, multipartidarismo e federalismo foram, por muito tempo, descritos como elementos potencialmente desorganizadores da autoridade governamental. Ainda assim, parte importante da literatura mostrou que esse arranjo não produziu, de modo necessário, paralisia decisória. Ao contrário, a experiência democrática brasileira revelou capacidade relevante de coordenação política, formação de maiorias legislativas e implementação de agendas governamentais \textcite{palermo2000_como_se} \textcite{limongi2017_a_crise}. O problema começa quando esse mesmo sistema, antes lido como funcional à governabilidade, passa a ser relido retrospectivamente como causa estrutural da ruptura presidencial de 2016.

Este paper parte da avaliação de que tal deslocamento interpretativo exige cautela. Tomar o impeachment de Dilma Rousseff como prova conclusiva da inviabilidade do presidencialismo de coalizão tende a reduzir um processo complexo a uma explicação institucional monocausal. Esse tipo de leitura perde de vista que a estabilidade presidencial depende não apenas de regras formais e instrumentos de coordenação, mas também da sustentação política das coalizões, da organização da competição partidária e da legitimidade das mediações representativas. Em outros termos, a crise não decorre automaticamente da arquitetura institucional; ela emerge quando mecanismos de cooperação deixam de operar sob condições políticas adversas.

A pergunta de pesquisa, portanto, é a seguinte: em que medida o impeachment de 2016 pode ser explicado menos como evidência da inviabilidade intrínseca do presidencialismo de coalizão e mais como resultado da interação entre erosão da base partidária de governo, dinâmica competitiva do sistema partidário e tensões de representação na democracia brasileira pós-1988? O objetivo geral do trabalho é analisar em que medida estabilidade e crise presidencial podem ser compreendidas a partir dessa interação, tendo o impeachment de 2016 como aplicação analítica central.

A hipótese defendida é que o impeachment de 2016 não demonstra, por si só, o fracasso estrutural do presidencialismo de coalizão. O arranjo institucional brasileiro mostrou capacidade de produzir governabilidade, mas essa capacidade depende da existência de coalizões partidárias operantes e de vínculos representativos minimamente estáveis. Quando o apoio legislativo se recompõe contra o Executivo e a representação partidária sofre erosão de legitimidade, conflitos antes administráveis podem ser convertidos em ruptura presidencial por vias formalmente democráticas. Nesse sentido, o impeachment deve ser interpretado como resultado de uma combinação entre crise de coalizão, reordenamento competitivo e deterioração das mediações políticas, não como simples efeito mecânico do desenho institucional \textcite{limongi2017_a_crise} \textcite{perez_lian2016_explaining_military}.

A relevância do tema é dupla. Do ponto de vista analítico, ele permite articular o bloco institucional da disciplina a pelo menos dois outros eixos decisivos: sistema partidário e representação, de um lado, e crise política e impeachment, de outro. Do ponto de vista substantivo, a questão ajuda a compreender por que a democracia brasileira pode combinar preservação das regras constitucionais com episódios agudos de destituição presidencial. Para desenvolver esse argumento, o paper organiza-se em quatro movimentos: primeiro, explicita problema, hipótese, objetivo e metodologia; segundo, reconstrói o debate teórico sobre governabilidade, presidencialismo de coalizão e representação; terceiro, aplica esse referencial ao caso de 2016; por fim, sintetiza as implicações analíticas da interpretação proposta.

\section{Problema de pesquisa e hipótese}
\label{sec:org66be366}
\subsection{Pergunta de pesquisa}
\label{sec:org5e7b481}
Em que medida o impeachment de 2016 pode ser explicado menos como evidência da inviabilidade intrínseca do presidencialismo de coalizão e mais como resultado da interação entre erosão da base partidária de governo, dinâmica competitiva do sistema partidário e tensões de representação na democracia brasileira pós-1988?

\subsection{Hipótese ou tese central}
\label{sec:org44ef1e9}
Este paper sustenta que o impeachment de 2016 não constitui, por si só, demonstração do fracasso estrutural do presidencialismo de coalizão. A literatura sobre a democracia brasileira indica que o arranjo institucional pós-1988 foi capaz de produzir coordenação decisória, disciplina legislativa e sustentação de governos, desde que lastreado por coalizões partidárias operantes \textcite{palermo2000_como_se} \textcite{limongi2017_a_crise}. A ruptura presidencial torna-se mais inteligível quando se considera a interação entre três processos: a erosão da base de apoio do Executivo no Legislativo, a recomposição competitiva do sistema partidário em chave adversa ao governo e o agravamento de tensões de representação que enfraquecem a legitimidade das mediações partidárias. Nessas condições, o impeachment funciona como mecanismo institucional de resolução de impasses, em consonância com a ideia de que, nas democracias latino-americanas, destituições legais podem desempenhar papel funcionalmente análogo ao de outros mecanismos históricos de remoção presidencial \textcite{perez_lian2016_explaining_military}.

\subsection{Objetivo geral}
\label{sec:org851e999}
O objetivo deste paper é analisar em que medida a estabilidade e a crise presidencial na democracia brasileira pós-1988 podem ser mais bem compreendidas pela interação entre capacidades institucionais de coordenação, recomposição das coalizões partidárias e tensões de representação no sistema partidário, tomando o impeachment de 2016 como aplicação analítica central.

\subsection{Metodologia}
\label{sec:org555ef6e}
O paper adota abordagem bibliográfica e argumentativa, articulando os textos-base da disciplina sobre instituições políticas, sistema partidário, representação, protesto e crise presidencial. Em vez de reconstruir exaustivamente os eventos de 2016, a estratégia consiste em mobilizar o caso como teste interpretativo de uma tese. O procedimento analítico combina, primeiro, reconstrução do debate sobre governabilidade e presidencialismo de coalizão; segundo, diálogo com interpretações sobre estrutura e inteligibilidade do sistema partidário; terceiro, incorporação de literatura sobre crise de representação, protestos e formas de engajamento político; e, por fim, aproximação com a literatura comparada sobre instabilidade presidencial \textcite{alonso2016_changing_repertoires} \textcite{singersd_tres_partidos} \textcite{alves2024_rage_within} \textcite{perez_lian2016_explaining_military}. O método, portanto, é qualitativo, bibliográfico e orientado por inferência argumentativa.

\section{Referencial teórico}
\label{sec:org3fc0d40}
\subsection{Governabilidade e presidencialismo de coalizão}
\label{sec:org61be9e6}
O debate sobre governabilidade no Brasil pós-1988 foi marcado, desde cedo, por diagnósticos pessimistas acerca da combinação entre presidencialismo e multipartidarismo. A imagem clássica de dispersão decisória pressupunha que o Executivo enfrentaria incentivos permanentes à instabilidade, dado o elevado número de partidos e a necessidade de negociar com um Congresso heterogêneo. Contra essa expectativa, parte relevante da literatura demonstrou que o sistema político brasileiro desenvolveu mecanismos robustos de coordenação. Em vez de um Executivo estruturalmente paralisado, observou-se capacidade de organizar maiorias, controlar a agenda e produzir decisão governamental \textcite{palermo2000_como_se}.

Em Palermo, o ponto decisivo está em deslocar a pergunta do plano normativo para o plano empírico: menos importante do que afirmar se o sistema concentra ou dispersa poder em abstrato é compreender como o governo efetivamente opera. Essa inflexão permite ver que a governabilidade brasileira não depende de homogeneidade partidária, mas de capacidade de coordenação institucionalmente ancorada. O presidencialismo de coalizão, nesse sentido, não é sinônimo de patologia, e sim forma concreta de organização do governo sob multipartidarismo. A questão analítica relevante passa a ser como se formam, se mantêm e se desfazem essas coalizões.

É precisamente nesse ponto que a intervenção de Limongi e Figueiredo adquire centralidade. Ao revisitar o debate em meio à crise do segundo governo Dilma Rousseff, os autores argumentam que o conceito de presidencialismo de coalizão foi progressivamente convertido em rótulo explicativo excessivamente elástico, passando a condensar sentidos normativos de cooptação, loteamento e corrupção que não derivam logicamente do conceito nem são suficientes para explicar a crise \textcite{limongi2017_a_crise}. A crítica é importante porque recoloca o foco sobre o funcionamento efetivo das instituições. Se o sistema foi capaz de produzir estabilidade ao longo de décadas, um episódio de ruptura não autoriza inferir, sem mediações, sua inviabilidade estrutural.

A formulação teórica que orienta este paper parte dessa revisão. O presidencialismo de coalizão deve ser entendido como mecanismo condicional de governabilidade: ele não elimina o conflito, mas o organiza. Sua eficácia depende da manutenção de incentivos à cooperação entre Executivo e partidos da base. Quando esses incentivos se alteram, a mesma estrutura que antes sustentava a coordenação pode deixar de proteger o governo. Assim, a análise da crise presidencial requer ir além da arquitetura formal e considerar a política de coalizões como processo dinâmico e reversível.

\subsection{Sistema partidário, competição e representação}
\label{sec:org97670a2}
A interpretação da crise presidencial exige, contudo, incorporar uma segunda dimensão: a do sistema partidário e de seus vínculos de representação. Uma visão recorrente da política brasileira sustenta que a multiplicidade de legendas produz um universo caótico, pouco inteligível e fracamente representativo. André Singer contesta esse diagnóstico ao argumentar que, sob a fragmentação aparente, a competição política brasileira apresentou organização discernível em torno de polos partidários com bases sociais e programáticas relativamente reconhecíveis \textcite{singersd_tres_partidos}. Isso não significa negar a fragmentação, mas relativizar a ideia de que ela dissolve toda mediação substantiva entre sociedade e política.

A contribuição de Singer é decisiva para este paper porque permite pensar o sistema partidário em chave relacional. O problema da governabilidade não é apenas o número de partidos, mas a forma como a competição se organiza e como os partidos conseguem representar clivagens sociais, interesses e identidades políticas. Se há algum grau de inteligibilidade na disputa partidária, então a crise de governo não pode ser lida simplesmente como produto automático da dispersão institucional. Ela deve ser situada no interior de um processo de reordenamento competitivo, no qual coalizões governamentais e oposicionistas disputam não só cargos e votos, mas também legitimidade representativa.

Esse ponto se torna ainda mais relevante quando se observa a crescente erosão da confiança nos partidos e nas mediações políticas tradicionais. A literatura sobre protestos e engajamento partidário sugere que a crise de representação brasileira não corresponde a um simples abandono da política, mas a uma reconfiguração de suas formas de expressão. Alonso e Mische mostram que junho de 2013 abriu um ciclo de protesto caracterizado por repertórios híbridos e forte rejeição aos partidos, repondo de maneira aguda o problema da relação entre mobilização social e política institucional \textcite{alonso2016_changing_repertoires}. Mais recentemente, Alves e coautores identificam a persistência de ativismo partidário intenso em contexto de desconfiança social, indicando que o vínculo com os partidos pode sobreviver sob modalidades negativas, polarizadas e movidas por antagonismo \textcite{alves2024_rage_within}.

Essas contribuições permitem qualificar a noção de crise de representação mobilizada neste paper. Não se trata apenas de déficit normativo de representatividade, mas de instabilidade das mediações por meio das quais partidos organizam apoios, traduzem demandas e estabilizam conflitos. Quando os partidos seguem relevantes, porém sob crescente polarização e deslegitimação recíproca, a competição tende a se tornar menos cooperativa. Nessa conjuntura, a governabilidade deixa de depender exclusivamente de instrumentos institucionais e passa a depender, de modo mais agudo, da capacidade de sustentar legitimidade política em ambiente hostil.

\subsection{Crise presidencial, impeachment e mediações institucionais}
\label{sec:org0302cdc}
A terceira dimensão do referencial teórico diz respeito à compreensão da ruptura presidencial dentro das regras democráticas. A literatura comparada de Pérez-Liñán e Polga-Hecimovich oferece um ponto de apoio importante ao sustentar que diferentes mecanismos de remoção presidencial podem funcionar como equivalentes funcionais, compartilhando causas gerais relacionadas à formação de amplas coalizões antagônicas ao Executivo \textcite{perez_lian2016_explaining_military}. Nas democracias recentes, isso não significa retorno dos golpes clássicos, mas maior ativação de instrumentos constitucionais de afastamento quando presidentes perdem suporte político decisivo.

Essa formulação é útil porque desloca o impeachment de uma leitura estritamente jurídica ou moral para uma leitura político-institucional. O impeachment é, ao mesmo tempo, mecanismo previsto na ordem constitucional e resultado de correlações de força que tornam sua ativação politicamente viável. Em outras palavras, a existência da regra não basta; é preciso que se forme uma coalizão capaz de convertê-la em desfecho. Isso aproxima a análise da crise presidencial do problema da governabilidade: presidentes caem não apenas porque existem instrumentos de destituição, mas porque perdem condições de coordenar apoios e bloquear adversários.

Embora o foco deste paper não seja federalismo nem judicialização, duas observações laterais ajudam a completar o quadro institucional. De um lado, a literatura sobre federalismo brasileiro mostra que a União construiu capacidades relevantes de coordenação e indução política ao longo da democracia recente, o que relativiza leituras de fragmentação estatal absoluta \textcite{oliveira2019_momentos_da}. De outro, estudos sobre o período posterior, especialmente sob Bolsonaro, indicam a expansão do protagonismo judicial em contextos de erosão política, sugerindo que crises de governabilidade tendem a redistribuir o peso relativo dos poderes e das arenas institucionais \textcite{vieira2022_supremocracia_e}. Essas observações reforçam a ideia de que o sistema brasileiro não é destituído de capacidades de coordenação; o problema é como tais capacidades se reorganizam quando a sustentação política do Executivo se rompe.

\subsection{Debate teórico e posicionamento analítico}
\label{sec:org33bffd0}
O debate teórico pode ser sintetizado em dois polos interpretativos. O primeiro lê a crise presidencial como confirmação de um defeito estrutural do arranjo institucional brasileiro: o presidencialismo de coalizão conteria, desde a origem, o germe de sua própria desestabilização. O segundo, mais compatível com Palermo e com Limongi e Figueiredo, argumenta que as instituições brasileiras revelaram capacidade duradoura de produzir governabilidade e que as rupturas exigem explicações mais situadas politicamente \textcite{palermo2000_como_se} \textcite{limongi2017_a_crise}.

Este paper se alinha ao segundo polo, mas com uma qualificação. Não basta afirmar que as instituições funcionavam e que, por isso, a crise deve ser atribuída a contingências externas. O argumento aqui é relacional: a governabilidade depende da articulação entre capacidades institucionais de coordenação, configuração da competição partidária e qualidade das mediações representativas. A crise emerge quando esses três planos deixam de convergir. Nessa chave, o impeachment de 2016 não refuta a existência de governabilidade prévia, mas mostra o caráter condicional e politicamente dependente do presidencialismo de coalizão.

Essa posição permite evitar dois reducionismos. O primeiro é o institucionalismo negativo que converte toda negociação de coalizão em sinônimo de degeneração do sistema. O segundo é o sociologismo difuso que explica a ruptura apenas por protestos, opinião pública ou polarização, sem considerar que tais processos precisam ser traduzidos em recomposição de maiorias institucionais para produzir impeachment. A contribuição analítica do paper consiste justamente em sustentar que a ruptura presidencial dentro da democracia brasileira pós-1988 é melhor compreendida como ponto de interseção entre instituições que organizam o conflito, partidos que o mediam e crises de representação que tornam essas mediações mais precárias.

\section{Análise}
\label{sec:org779ee01}
\subsection{Recorte empírico ou problema analisado}
\label{sec:orgb88edc8}
O recorte empírico considerado neste paper é o impeachment de 2016 como momento de conversão de uma crise de governabilidade em ruptura presidencial dentro das regras democráticas. O foco não está na reconstituição cronológica do episódio, nem na avaliação jurídica de seus fundamentos, mas no uso do caso como aplicação analítica para testar a hipótese formulada. A questão central é saber por que um sistema que havia produzido estabilidade e coordenação por largo período deixou de proteger o mandato presidencial.

Tomado isoladamente, o impeachment pode sugerir a imagem de um arranjo institucional inviável. No entanto, quando situado no percurso da democracia pós-1988, ele aparece menos como evidência suficiente de falha estrutural e mais como desfecho de uma deterioração específica das condições políticas de governabilidade. A literatura de Palermo e de Limongi e Figueiredo oferece o ponto de partida para essa leitura: se o sistema mostrou capacidade reiterada de organizar decisões e formar maiorias, então a explicação da ruptura precisa incidir sobre a transformação dessas maiorias e sobre a erosão das mediações que lhes davam sustentação \textcite{palermo2000_como_se} \textcite{limongi2017_a_crise}.

Nessa perspectiva, o impeachment de 2016 é analiticamente relevante porque combina três processos. Primeiro, a perda progressiva de coordenação da coalizão governista no interior do Legislativo. Segundo, a reconfiguração da competição partidária em ambiente de forte antagonismo e incentivo à deserção. Terceiro, a ampliação de tensões de representação que reduziram a capacidade dos partidos governistas de estabilizar social e simbolicamente o conflito. O episódio torna-se, assim, um caso limite para observar quando a engrenagem do presidencialismo de coalizão deixa de operar como mecanismo de produção de governabilidade e passa a funcionar como arena de recomposição adversa do poder.

\subsection{Aplicação dos conceitos ao caso}
\label{sec:orgeec92fc}
A primeira implicação do referencial teórico é rejeitar a ideia de que o presidencialismo de coalizão colapsa por definição. Se isso fosse verdadeiro, a democracia brasileira pós-1988 teria apresentado padrão contínuo de paralisia ou de rupturas frequentes, o que não corresponde ao argumento da literatura mobilizada. O que o caso de 2016 sugere, em vez disso, é que a governabilidade depende de uma base partidária efetiva, não apenas formal. Quando a coalizão se torna incapaz de oferecer previsibilidade legislativa e coordenação estratégica ao Executivo, a presidência perde sua principal defesa institucional.

Esse ponto recoloca o impeachment no terreno da política de coalizões. Uma vez erodido o apoio partidário, o problema deixa de ser apenas governar com dificuldade; passa a ser sobreviver institucionalmente. A tese de Pérez-Liñán e Polga-Hecimovich ajuda a esclarecer esse deslocamento: presidentes são removidos quando se formam coalizões suficientemente amplas e motivadas para ativar mecanismos de destituição \textcite{perez_lian2016_explaining_military}. No Brasil de 2016, essa chave sugere que o impeachment dependeu menos de uma simples falha do desenho constitucional e mais da emergência de uma maioria antagônica capaz de transformar conflito político em desfecho institucional.

Entretanto, a formação dessa maioria não pode ser entendida apenas como rearranjo parlamentar autônomo. Ela esteve vinculada à dinâmica mais ampla do sistema partidário. A contribuição de Singer é relevante porque permite enxergar o conflito não como produto aleatório de um sistema pulverizado, mas como expressão de uma disputa estruturada entre polos partidários e sociais \textcite{singersd_tres_partidos}. O impeachment, nesse sentido, deve ser lido também como momento de reordenação da competição. Partidos e lideranças avaliaram custos e benefícios de manter ou abandonar a coalizão governista em um contexto no qual a continuidade do apoio ao Executivo já não maximizava suas expectativas de poder.

A erosão da governabilidade, portanto, não foi mero problema administrativo. Ela foi inseparável de uma crise de representação. Os protestos iniciados em 2013 contribuíram para deslocar o centro de gravidade da disputa política, ampliando a rejeição aos partidos e enfraquecendo a capacidade das mediações institucionais de absorver descontentamentos de forma rotinizada \textcite{alonso2016_changing_repertoires}. Isso não significa afirmar que junho levou linearmente ao impeachment, mas reconhecer que o ciclo de protesto reabriu a relação entre rua, opinião pública e legitimidade dos partidos. Em tal ambiente, a defesa do governo tornou-se mais custosa para aliados, enquanto sua deposição se tornou mais politicamente inteligível e menos onerosa.

A análise de Rocha acrescenta uma dimensão organizativa importante a esse quadro. Ao mostrar a formação de um contrapúblico ultraliberal que teve papel relevante nos protestos pró-impeachment, a autora sugere que a crise não se limitou à perda passiva de apoio ao governo, mas envolveu também a construção ativa de campos militantes e repertórios de mobilização orientados para sua destituição \textcite{rocha2019_imposto_e}. Isso reforça a ideia de que tensões de representação não operam apenas como ausência de mediação; elas podem favorecer novas formas de articulação política que pressionam por soluções institucionais radicais dentro da legalidade democrática.

Ao mesmo tempo, o caso de 2016 qualifica a leitura segundo a qual a fragmentação partidária, isoladamente, explicaria a ruptura. Se a fragmentação fosse causa suficiente, a instabilidade presidencial seria contínua. O que importa analiticamente é a fragmentação sob condições específicas de competição e deslegitimação. Em cenário ordinário, o presidencialismo de coalizão pode administrar heterogeneidade partidária. Em cenário extraordinário, quando a coalizão governista se desmancha, a oposição percebe oportunidade estratégica e a representação partidária perde legitimidade pública, a mesma heterogeneidade pode acelerar a recomposição de maiorias contra o Executivo. O problema, portanto, não é o multipartidarismo em abstrato, mas sua interação com ciclos de conflito e com incentivos de curto prazo à deserção.

A literatura recente sobre engajamento partidário sob polarização ajuda a compreender por que a crise de representação não produz necessariamente despolitização. Alves e coautores mostram que, em contextos de forte antagonismo, o ativismo pode ser intensificado por afetos negativos e identidades reativas \textcite{alves2024_rage_within}. Essa observação é útil para pensar retrospectivamente o ambiente em que o impeachment se tornou possível: a deterioração da confiança nos partidos não impediu mobilização; ao contrário, favoreceu formas de engajamento marcadas por rejeição intensa ao adversário. Com isso, a competição partidária se tornou mais conflitiva e menos propensa a compromissos cooperativos, enfraquecendo as bases informais da governabilidade.

Por fim, a análise do caso mostra que a ruptura presidencial não implica colapso da ordem democrática, mas transformação da forma de administrar conflitos. Essa é justamente a força da abordagem comparada de Pérez-Liñán: em democracias latino-americanas, o afastamento presidencial pode ocorrer por meios constitucionais quando a coalizão governante perde sua capacidade de proteção \textcite{perez_lian2016_explaining_military}. No Brasil, o impeachment de 2016 expressa esse padrão. Houve ruptura de mandato, mas não suspensão da ordem constitucional. Isso torna a explicação mais exigente: não basta falar em crise; é preciso explicar como uma crise política se traduziu em destituição por meio de instituições regulares.

\subsection{Discussão crítica}
\label{sec:orgecc7a3c}
A principal consequência analítica do argumento desenvolvido é que o impeachment de 2016 não deve ser interpretado como refutação simples da literatura sobre governabilidade no Brasil. Ao contrário, ele confirma indiretamente um de seus pressupostos centrais: governos dependem de coalizões partidárias operantes. O erro de certa leitura pós-crise consistiu em transformar a necessidade de negociação em prova de degeneração estrutural. Essa operação apaga a distinção entre mecanismos de coordenação e os contextos políticos nos quais tais mecanismos deixam de funcionar \textcite{limongi2017_a_crise}.

Ao mesmo tempo, seria insuficiente responder à crítica apenas reafirmando a robustez das instituições. O caso de 2016 mostra que a estabilidade produzida pelo presidencialismo de coalizão é politicamente condicionada. A governabilidade não é propriedade automática do desenho institucional; ela resulta da convergência entre regras, incentivos partidários e legitimidade representativa. Quando uma dessas dimensões se deteriora, o sistema pode absorver o conflito. Quando as três se desalinham simultaneamente, a ruptura presidencial se torna uma possibilidade concreta.

Essa formulação permite reinterpretar o sistema partidário brasileiro de forma menos esquemática. Nem caos sem representação, nem estrutura perfeitamente estável. A competição partidária brasileira mostrou, por longo período, capacidade de ordenar disputas e sustentar governos, como sugere Singer \textcite{singersd_tres_partidos}. Mas essa ordenação era dependente de mediações que se fragilizaram sob polarização, protesto anti-institucional e recomposição estratégica das elites partidárias. Em outras palavras, havia inteligibilidade competitiva, mas não imunidade à crise.

Também é importante notar que a explicação proposta evita a ilusão de causalidade única. O impeachment não decorre apenas da crise de representação, nem apenas da deserção parlamentar, nem apenas da forma institucional do presidencialismo. Sua inteligibilidade está na combinação desses elementos. A contribuição específica do presidencialismo de coalizão para o argumento não é ser a causa da crise, mas ser a forma normal de organização do governo cuja sustentação se rompeu. A partir dessa ruptura, o mecanismo constitucional do impeachment tornou-se politicamente utilizável.

Em sentido mais amplo, o caso de 2016 sugere uma lição sobre a democracia brasileira pós-1988. Sua resiliência institucional não elimina vulnerabilidades; ela apenas desloca a forma pela qual a instabilidade se manifesta. Em vez de golpes militares, predominam conflitos resolvidos no interior das regras, ainda que com alta intensidade e forte contestação política. Isso não reduz a gravidade da ruptura, mas exige linguagem conceitual mais precisa. O problema não é a inexistência de instituições, e sim a precarização das mediações que permitem a essas instituições produzir cooperação e legitimidade ao mesmo tempo.

\section{Conclusão}
\label{sec:org368fc47}
Em síntese, o paper argumentou que o impeachment de 2016 é mais bem compreendido não como prova conclusiva da inviabilidade intrínseca do presidencialismo de coalizão, mas como resultado da interação entre erosão da base partidária de governo, reconfiguração competitiva do sistema partidário e agravamento das tensões de representação na democracia brasileira pós-1988. A pergunta de pesquisa foi respondida por meio de um argumento relacional: a ruptura presidencial ocorre quando as capacidades institucionais de coordenação deixam de estar sustentadas por coalizões operantes e por mediações representativas minimamente estáveis.

A análise permitiu reafirmar a hipótese central. A literatura de Palermo e de Limongi e Figueiredo mostra que o arranjo institucional brasileiro foi capaz de produzir governabilidade, o que impede tomar um episódio de ruptura como evidência suficiente de fracasso estrutural \textcite{palermo2000_como_se} \textcite{limongi2017_a_crise}. Ao mesmo tempo, o caso de 2016 exige qualificar essa conclusão: a estabilidade presidencial é condicional e depende de convergência entre instituições, partidos e representação. Quando o Executivo perde proteção legislativa, a competição partidária se reorganiza em chave adversa e a legitimidade das mediações políticas se enfraquece, o impeachment torna-se uma via institucional plausível de resolução do conflito, como sugere a literatura comparada sobre instabilidade presidencial \textcite{perez_lian2016_explaining_military}.

A principal contribuição analítica do paper consiste em recusar tanto a explicação monocausal centrada no desenho institucional quanto leituras difusas que dissolvem a crise em mera insatisfação social. O presidencialismo de coalizão permanece útil como conceito, desde que entendido como mecanismo condicional de governabilidade e não como rótulo totalizante da crise. Nessa chave, o impeachment de 2016 ilumina menos um defeito originário da democracia brasileira do que a fragilidade das mediações políticas que tornam suas instituições operantes.

\printbibliography
\end{document}